\section{Datasets Used}
\label{sec:datasets}

\subsection{The current dataset}

The working dataset is \texttt{unitime-out/0gen\_all}, a complete set of the sixteen Data Exchange
documents for a single academic session. Every file carries the term attribute \texttt{Spr12} and
the year \texttt{2026}, giving the Data Exchange term identifier \texttt{Spr122026}. The session
runs from 19 January 2026 to 22 May 2026, with classes ending on 30 April and the examination
period beginning on 4 May. The weekly grid is seven days of ten hourly slots from 08:30 to 18:30,
with the underlying solver resolution being five-minute slots.

Table~\ref{tab:dataset} gives the inventory, recomputed directly from the files rather than quoted
from documentation.

\begin{table}[H]
\centering
\caption{Inventory of \texttt{unitime-out/0gen\_all} (term \texttt{Spr122026}).}
\label{tab:dataset}
\footnotesize
\begin{tabular}{@{}rL{5.0cm}rL{6.6cm}@{}}
\toprule
\textbf{\#} & \textbf{File} & \textbf{Bytes} & \textbf{Contents} \\
\midrule
1  & \texttt{1sessionSetup.xml}            &  12{,}486 & session, departments, subject areas, date and time patterns \\
2  & \texttt{2academicArea.xml}            &     496 & 2 areas (B.Tech, M.Tech) \\
3  & \texttt{3academicClassification.xml}  &     786 & 5 year codes (FY, SY, TY, BT, MT) \\
4  & \texttt{4Major.xml}                   &   1{,}832 & 15 majors \\
5  & \texttt{5Minor.xml}                   &     374 & empty by design; Taasika models minors as majors \\
6  & \texttt{6studentGroup.xml}            &  25{,}485 & \textbf{242 groups} (43 classes + 199 batches) \\
7  & \texttt{7buildingRoomImport.xml}      &  21{,}604 & \textbf{6 buildings, 42 rooms} \\
8  & \texttt{8roomSharing.xml}             &   9{,}916 & 45 sharing blocks; weekends unavailable \\
9  & \texttt{9travelTimes.xml}             &   2{,}722 & 7 inter-building travel entries \\
10 & \texttt{10staff.xml}                  &  19{,}618 & \textbf{115 instructors} \\
11 & \texttt{11courseCatalog.xml}          &  44{,}081 & \textbf{158 courses} \\
12 & \texttt{12courseOffering.xml}         & 100{,}340 & \textbf{56 offerings, 84 subparts, 291 classes} \\
13 & \texttt{13preferences.xml}            &  95{,}146 & 362 preference blocks \\
14 & \texttt{14studentInfo.xml}            & 657{,}518 & \textbf{1{,}130 students}, 5{,}767 group references \\
15 & \texttt{15studentRequest.xml}         & 495{,}638 & course requests for 1{,}130 students \\
16 & \texttt{16studentenrollments.xml}     & 661{,}230 & \textbf{9{,}148 section enrolments} \\
\bottomrule
\end{tabular}
\end{table}

\subsection{Source snapshot}

Determining which Taasika snapshot produced this dataset required evidence rather than trust. Four
of the files carry the attribute \texttt{created="Generated from Taasika snapshot 240"}, but that
string is hard-coded in three of the generator scripts and is not updated when a different snapshot
is selected, so it cannot be used as provenance. The snapshot was therefore identified by counting.

\begin{table}[H]
\centering
\caption{Identifying the source snapshot by counting entities.}
\label{tab:snapshot}
\footnotesize
\begin{tabular}{@{}L{4.6cm}rrrl@{}}
\toprule
\textbf{Quantity} & \textbf{Snap 240} & \textbf{Snap 242} & \textbf{Snap 263} & \textbf{\texttt{0gen\_all}} \\
\midrule
Classes + batches = groups & $43{+}165={}$208 & $43{+}199={}$\textbf{242} & $32{+}128={}$160 & \textbf{242} \\
Teachers                   & 113 & \textbf{115} & 67 & \textbf{115} \\
Rooms                      & 49  & 49  & 29 & 42 (see below) \\
\bottomrule
\end{tabular}
\end{table}

The dataset was generated from snapshot \textbf{242}, named \texttt{published-v4-26jan26}, in
\texttt{taasika-db/taasika3.sql}. Snapshot 240 is the Semester~VI dataset and is what
\texttt{unitime-out/README.md} still documents; snapshot 263 is a later snapshot for the following
academic year that was integrated into the loader but is not the source of these files. The stale
provenance string is recorded as a known defect in Section~\ref{sec:limitations}.

The room count differs from the snapshot because the pipeline no longer emits the virtual
``Online'' building: 49 rooms in the snapshot less seven online rooms gives the 42 that are
imported. The six buildings are the Academic Complex (14 rooms), the Computer Engineering and IT
department (15), the ENTC Extension (5), the New CSE Building (2), the Bhau Institute (2), and a
pool of four laboratories marked unavailable, which is retained deliberately so that any use of
them remains visible in the evaluation report.

\subsection{Student population}

Institutional MIS data was not available, so the student body is simulated from headcounts
confirmed with the department and recorded in \texttt{scripts/intake.py}.
Table~\ref{tab:students} gives the composition. Identifiers follow the pattern \texttt{61yy03aaa},
where \texttt{yy} encodes the year of study. No real student data appears anywhere in the
repository or in this report.

\begin{table}[H]
\centering
\caption{Simulated student population of the current dataset.}
\label{tab:students}
\footnotesize
\begin{tabular}{@{}lrrrL{6.2cm}@{}}
\toprule
\textbf{Year} & \textbf{CSE} & \textbf{AI\,\&\,ML} & \textbf{Total} & \textbf{Structure} \\
\midrule
First year (FY)   & 300 & 100 & 400   & single cohort, reserved blocks \\
Second year (SY)  & 200 & 0   & 200   & two lecture divisions; direct second-year entrants join here \\
Third year (TY)   & 200 & 0   & 200   & two divisions, departmental electives \\
Final year (BT)   & 200 & 0   & 200   & two divisions, one departmental elective each \\
M.Tech (MT)       & \multicolumn{2}{c}{130} & 130 & specialisation streams \\
\midrule
\textbf{Total}    &     &     & \textbf{1{,}130} & all with enrolments \\
\bottomrule
\end{tabular}
\end{table}

Minor-degree and open-elective provision is modelled as room blocks rather than as enrolments,
because those courses are taught by other departments: 600 seats are reserved for minor-degree
teaching across the second, third and final years, and 200 for the second-year open elective.

Two known imperfections in the simulation are recorded rather than hidden. Two hundred first-year
students could not be placed in \texttt{CS PPS-FY Lec} because the declared section limits do not
cover the cohort, and two M.Tech students overflow the \texttt{PSEC2-EAI} laboratory. The section
filling routine assigns students to a preferred section up to its declared limit and then spills
into later sections of the same course and type, which is also why a small number of laboratory
batches contain students drawn from two different lecture divisions.

\subsection{Other instances used in the experiments}

Several earlier instances appear in Section~\ref{sec:experiments}, and their counts differ. The
term labels are \emph{not} different semesters: \texttt{Spr6}, \texttt{Spr10}, \texttt{Spr11} and
\texttt{Spr12} are successive UniTime academic sessions created for Spring 2026, each a clean
re-import of an evolving dataset. Table~\ref{tab:instances} lists them so that no figure in this
report appears without its denominator.

\begin{table}[H]
\centering
\caption{Instances referenced in this report.}
\label{tab:instances}
\footnotesize
\begin{tabular}{@{}L{4.4cm}lrrrL{4.0cm}@{}}
\toprule
\textbf{Instance} & \textbf{Session} & \textbf{Classes} & \textbf{Students} & \textbf{Rooms} & \textbf{Role} \\
\midrule
\texttt{rl/instances/spr2026.xml}        & Spr6  & 262 & 990   & 44 & solver export used for all headless experiments \\
\texttt{spr2026\_10m\_w5.xml}            & Spr6  & 262 & 990   & 44 & already conflict-free; bake-off instance \\
\texttt{spr2026\_10\_LE\_SS\_Enrolled}   & Spr10 & 290 & 1{,}130 & 41 & first pinned-enrolment run \\
\texttt{spr2026\_11\_all\_flags.csv}     & Spr11 & 291 & 1{,}130 & 42 & final solution, scored in Section~\ref{sec:experiments} \\
\texttt{unitime-out/0gen\_all}           & Spr12 & 291 & 1{,}130 & 42 & current input dataset \\
\bottomrule
\end{tabular}
\end{table}

The 990 figure recurring in the Spr6 rows is the number of students that the UniTime solver export
carried enrolments for, out of a simulated roll of 1{,}290 in the dataset of that period. The gap
is not attrition; it is the difference between the students the solver loaded and the students that
existed in the source files, and it is the origin of the conflict-rate reconciliation in
Section~\ref{sec:experiments}.
